\chapter{Procedimentos de Teste e Validação do Projeto}
\label{cap:testes}

% Nesta seção, devem ser indicados como os diversos módulos do projeto serão testados e validados individualmente e em conjunto.
% Deverão ser estabelecidos os critérios de aceite da solução em termos de desempenho e cumprimento de seus aspectos funcionais.

\section{Procedimentos de Teste}

A validação da solução será conduzida de forma modular e experimental. Inicialmente, cada módulo do \textit{pipeline} será verificado funcionalmente:
\begin{itemize}
    \item \textbf{Ingestão e normalização}: verificação do processamento de arquivos válidos e tratamento de entradas inconsistentes, garantindo o isolamento da unidade fundamental de análise, que é o \textbf{método ou função} modificada no \textit{commit} ou \textit{pull request};
    \item \textbf{Análise estrutural e extração}: verificação da geração das representações sintáticas básicas (como nós da AST e relações de chamadas) a partir de exemplos controlados em Java (e C\# na etapa estendida);
    \item \textbf{Preparação e partição dos dados}: conferência da integridade dos rótulos e garantia de que instâncias de um mesmo repositório não sejam compartilhadas entre os conjuntos de treinamento e teste, mitigando o risco de vazamento de dados \cite{ding2024primevul}.
\end{itemize}

A avaliação do módulo de classificação será realizada em um conjunto de teste independente, registrando-se a matriz de confusão e as métricas convencionais de desempenho: precisão, revocação e medida F1. 

Para contextualizar os resultados obtidos, a solução será comparada a uma \textbf{linha de base (baseline) de referência}, constituída pela execução de regras padrão de detecção para as quatro CWEs em uma ferramenta SAST aberta consagrada (como o Semgrep). O objetivo dessa comparação experimental é avaliar se a adição do classificador de \textit{Machine Learning} sobre as características extraídas proporciona uma tendência de redução de falsos positivos frente aos alertas brutos da análise estática, buscando manter um equilíbrio adequado com a taxa de revocação para não descartar falhas de segurança reais.

Por fim, a integração entre os módulos será testada em fluxo contínuo, assegurando que o código fornecido como entrada percorra as etapas de processamento e resulte na emissão estruturada do relatório de alertas ao usuário.

\newpage

\section{Critérios de Aceite e Desempenho}

Os critérios de aceite da solução são apresentados na Tabela~\ref{tab:criterios_aceite}.

\begin{longtable}{p{0.25\textwidth} p{0.68\textwidth}}
    \caption{Critérios de aceite da solução.}
    \label{tab:criterios_aceite} \\
    \toprule
    \textbf{Aspecto} & \textbf{Critério de aceite} \\
    \midrule
    \endfirsthead
    \toprule
    \textbf{Aspecto} & \textbf{Critério de aceite} \\
    \midrule
    \endhead
    Entrada e normalização & Processar arquivos de código em Java (foco do MVP) e C\#, isolando as funções modificadas e emitindo mensagens claras em caso de erros de entrada. \\
    \addlinespace
    Extração de \\características & Gerar as estruturas sintáticas esperadas e extrair atributos representativos para as funções analisadas, preservando o mapeamento de arquivo e linha. \\
    \addlinespace
    Gestão de dados & Separar os conjuntos de treino e teste por repositório ou projeto, assegurando avaliação não enviesada. \\
    \addlinespace
    Classificação e baseline & Apresentar as métricas de precisão, revocação e F1 no conjunto de teste. O desempenho será comparado à linha de base SAST, avaliando o potencial de redução de falsos positivos e a capacidade de retenção de vulnerabilidades reais. \\
    \addlinespace
    Relatório de saída & Emitir os alertas identificando a CWE associada, o arquivo, a linha aproximada e o nível de confiança da predição. \\
    \addlinespace
    Integração & Executar o fluxo completo do sistema, desde a leitura do código-fonte até a geração do relatório, sem falhas de comunicação entre os blocos. \\
    \addlinespace
    Reprodutibilidade & Garantir que o ambiente de execução e os experimentos possam ser reproduzidos com os mesmos conjuntos de dados e configurações. \\
    \bottomrule
\end{longtable}
